% Blueprint for the formalization of
%   "Integer Cantor Sets: Arithmetic Combinatorial Properties"
%   Burgin, Fragkos, Lacey, Mena, Reguera.
%
% Conventions used below.
%   \lean{Name}   the Lean declaration implementing the statement
%   \leanok       on a statement: the Lean statement exists and typechecks
%                 on a proof:     the Lean proof is complete (no sorry)
%   \uses{...}    dependencies, used to build the dependency graph
%   \notready     the statement depends on material not yet in Mathlib

\chapter{Scope and strategy}

\section{What this blueprint covers}

The paper establishes four groups of results about \emph{integer Cantor sets}
$\Cbd$, the non-negative integers all of whose base $b$ digits lie in a fixed
set $D$: (a) intersective and van der Corput properties, (b) polynomial
uniform distribution of the pair $(k, \sd(k))$, (c) mean ergodic theorems, and
(d) metric Poissonian pair correlation.

These four groups differ enormously in formalization cost. The purpose of this
chapter is to say so explicitly, before any mathematics, so that effort is
spent where it will actually close.

\begin{itemize}
  \item \textbf{Chapters \ref{chap:structure} and \ref{chap:vdc} are
    self-contained.} They need only \texttt{Nat.digits}, finite sums, and
    elementary real analysis. Everything there should close without new
    Mathlib infrastructure. This is where the formalization starts.
  \item \textbf{Chapters \ref{chap:expsum} and \ref{chap:pudu} are hard but
    unblocked.} They need limits of infinite products of complex numbers,
    Weyl's criterion, and a two-variable polynomial induction. All of the
    ingredients exist in Mathlib; the work is long but ordinary. Both are now
    complete.
  \item \textbf{Chapter \ref{chap:residue} is unblocked and moderate.} It is
    finite Fourier analysis on $\Z_q$ plus the results of Chapter
    \ref{chap:expsum}.
  \item \textbf{Chapter \ref{chap:recurrence} is blocked} on the Furstenberg
    correspondence principle, which is not in Mathlib.
  \item \textbf{Chapter \ref{chap:ergodic} is blocked} on the joint spectral
    measure of a pair of commuting unitary operators, which is not in Mathlib.
  \item \textbf{Chapter \ref{chap:paircorr} is blocked} on the theorem of
    Aistleitner, Larcher and Lewko, which is a research paper in its own right
    and is not a realistic formalization target. It is recorded as a
    hypothesis.
\end{itemize}

\section{Mathlib inventory}

The following are available and will be used heavily.

\begin{itemize}
  \item \texttt{Nat.digits b n : List Nat}, the little-endian base $b$ digit
    list, together with \texttt{Nat.ofDigits}, \texttt{Nat.digits\_lt\_base},
    \texttt{Nat.digits\_ne\_nil\_iff\_ne\_zero},
    \texttt{Nat.getLast\_digit\_ne\_zero} and
    \texttt{Nat.digits\_ofDigits} (a round trip valid exactly when the digits
    are bounded by $b$ and the leading digit is non-zero).
  \item \texttt{Nat.modEq\_digits\_sum}, which gives $b-1 \mid k - \sd(k)$
    directly.
  \item \texttt{Finset.sort}, \texttt{Finset.orderIsoOfFin} for the increasing
    enumeration of a finite set of naturals.
  \item \texttt{Complex.exp}, \texttt{Real.cos}, \texttt{Finset.sum}, and the
    standard limit API.
  \item \texttt{Subadditive} and \texttt{Subadditive.tendsto\_lim}, Fekete's
    lemma, which gives convergence of $I(H,N)$ in
    Lemma \ref{lem:intersective-converges}.
\end{itemize}

The following are \emph{not} available and are the real obstructions.

\begin{itemize}
  \item The Furstenberg correspondence principle.
  \item Joint spectral measures / the spectral theorem for commuting unitaries
    on a Hilbert space. Mathlib has the continuous functional calculus and the
    spectral theorem in finite dimensions and for compact operators, but not
    the projection-valued measure on $\T^2$ that Chapter \ref{chap:ergodic}
    needs.
  \item Upper Banach density, and the theory of IP and IP$^*$ sets.
  \item Metric Poissonian pair correlation.
\end{itemize}

\section{A remark on the standing assumption}

The paper's standing assumption is that $b \geq 3$, $D \subseteq
\{0,\ldots,b-1\}$, $0 \in D$, and $2 \leq \abs{D} < b$. The hypothesis $0 \in
D$ is doing real work, and the formalization makes the reason unusually sharp.

The increasing enumeration $\enum$ of $\Cbd$ is defined (Definition
\ref{def:enum}) by reading the base $\abs{D}$ digits of $n$ and substituting
the $i$-th smallest element of $D$ for the digit $i$. For this to be the
increasing enumeration one needs \texttt{Nat.digits\_ofDigits} to apply, and
that lemma requires the \emph{leading} digit of the output to be non-zero.
The leading base $\abs{D}$ digit of $n$ is non-zero, so the substituted digit
is $d_i$ for some $i \geq 1$; this is non-zero precisely because $d_0 = 0$,
i.e. precisely because $0 \in D$. Without $0 \in D$ the substitution is not
even digit-preserving, and Lemma \ref{lem:enum-block} fails: for $b = 3$ and
$D = \{1,2\}$ one has $\enum(2) = 4$, whereas the formula would predict
$1 + 2\cdot 3 = 7$.

\chapter{Cantor sets: definitions and structure}
\label{chap:structure}

Throughout, $b$ is a natural number and $D$ a finite set of natural numbers.
We write $\delta = \abs{D}$ for the number of allowed digits.

\section{Definitions}

\begin{definition}[Integer Cantor set]
  \label{def:cantorSet}
  \lean{IntegerCantorSet.cantorSet}
  \leanok
  The \emph{integer Cantor set} with base $b$ and digit set $D$ is
  \[
    \Cbd = \{\, n \in \N \;:\; \text{every base } b \text{ digit of } n
      \text{ lies in } D \,\}.
  \]
  In Lean this is \texttt{\{n | $\forall$ d $\in$ Nat.digits b n, d $\in$ D\}}.
  Note that $0 \in \Cbd$ for every $D$, since \texttt{Nat.digits b 0 = []}.
\end{definition}

\begin{definition}[Standing assumption]
  \label{def:standing}
  \lean{IntegerCantorSet.Standing}
  \leanok
  \uses{def:cantorSet}
  A pair $(b, D)$ satisfies the \emph{standing assumption} if
  \[
    b \geq 3, \qquad D \subseteq \{0, 1, \ldots, b-1\}, \qquad 0 \in D,
    \qquad 2 \leq \abs{D} < b .
  \]
\end{definition}

\begin{definition}[Sum of digits]
  \label{def:sumDigits}
  \lean{IntegerCantorSet.sumDigits}
  \leanok
  $\sd(n)$ is the sum of the base $b$ digits of $n$, that is
  \texttt{(Nat.digits b n).sum}.
\end{definition}

\begin{definition}[Step size]
  \label{def:stepSize}
  \lean{IntegerCantorSet.stepSize}
  \leanok
  $s = s(D)$ is the greatest common divisor of $D$. Under the standing
  assumption this is the largest $s$ such that $D$ is contained in an
  arithmetic progression of step $s$, because $0 \in D$.
\end{definition}

\begin{definition}[The $i$-th digit]
  \label{def:digitOf}
  \lean{IntegerCantorSet.digitOf}
  \leanok
  $d_i$ is the $i$-th smallest element of $D$, with the junk value $0$ for
  $i \geq \delta$. In Lean, \texttt{(D.sort (· $\leq$ ·)).getD i 0}.
\end{definition}

\begin{definition}[Increasing enumeration]
  \label{def:enum}
  \lean{IntegerCantorSet.enum}
  \leanok
  \uses{def:digitOf}
  For $n \in \N$ write $n = \sum_c n_c \delta^c$ in base $\delta$. Then
  \[
    \enum(n) = \sum_c d_{n_c} b^c ,
  \]
  in Lean \texttt{Nat.ofDigits b ((Nat.digits D.card n).map (digitOf D))}.
\end{definition}

\section{Elementary properties of the digit substitution}

\begin{lemma}
  \label{lem:digitOf-zero}
  \lean{IntegerCantorSet.digitOf_zero}
  \leanok
  \uses{def:digitOf, def:standing}
  If $0 \in D$ and $D$ is non-empty then $d_0 = 0$.
\end{lemma}
\begin{proof}
  \uses{lem:digitOf-zero}
  \leanok
  $D.\mathrm{sort}$ is sorted in increasing order, so its $0$-th entry is
  $\min D$ (\texttt{Finset.sorted\_zero\_eq\_min'}), and $0$ is the least
  element of any set of naturals containing $0$.
\end{proof}

\begin{lemma}
  \label{lem:digitOf-mem}
  \lean{IntegerCantorSet.digitOf_mem}
  \leanok
  \uses{def:digitOf}
  For $i < \delta$ we have $d_i \in D$.
\end{lemma}
\begin{proof}
  \uses{lem:digitOf-mem}
  \leanok
  The sorted list of $D$ has length $\delta$ and its entries are exactly the
  members of $D$.
\end{proof}

\begin{lemma}
  \label{lem:digitOf-strictMono}
  \lean{IntegerCantorSet.digitOf_strictMonoOn}
  \leanok
  \uses{def:digitOf}
  $i \mapsto d_i$ is strictly increasing on $\{0, \ldots, \delta - 1\}$.
\end{lemma}
\begin{proof}
  \uses{lem:digitOf-strictMono}
  \leanok
  \texttt{Finset.sort} produces a list that is sorted and has no duplicates,
  hence strictly sorted (\texttt{Finset.sortedLT\_sort}); injectivity follows
  from strict monotonicity by trichotomy.
\end{proof}

\begin{lemma}
  \label{lem:digitOf-ne-zero}
  \lean{IntegerCantorSet.digitOf_ne_zero}
  \leanok
  \uses{lem:digitOf-zero, lem:digitOf-strictMono}
  Assume the standing assumption. For $1 \leq i < \delta$ we have $d_i \neq 0$.
\end{lemma}
\begin{proof}
  \uses{lem:digitOf-ne-zero}
  \leanok
  By Lemma \ref{lem:digitOf-zero}, $d_0 = 0$, and by Lemma
  \ref{lem:digitOf-strictMono}, $d_i > d_0$ for $i \geq 1$.
\end{proof}

\section{The enumeration theorem}

This is Proposition 2.5 of the paper. It is the technical heart of
everything that follows, and it is where $0 \in D$ is used.

\begin{lemma}[Single-digit recursion]
  \label{lem:enum-recursion}
  \lean{IntegerCantorSet.enum_mul_add,
    IntegerCantorSet.enum_of_lt_card}
  \leanok
  \uses{def:enum, lem:digitOf-zero}
  Assume the standing assumption. For $j < \delta$ one has
  $\enum(j) = d_j$, and for every $m$,
  \[
    \enum(\delta m + j) = b\,\enum(m) + d_j .
  \]
\end{lemma}
\begin{proof}
  \uses{lem:enum-recursion}
  \leanok
  For $m = 0$ this is the first claim, which holds because
  $\dgt_\delta(j) = [j]$ for $0 < j < \delta$ and because $d_0 = 0$
  (Lemma \ref{lem:digitOf-zero}) when $j = 0$.  For $m > 0$,
  $\delta m + j > 0$, so \texttt{Nat.digits\_def'} gives
  $\dgt_\delta(\delta m + j) = j :: \dgt_\delta(m)$; substituting and
  expanding \texttt{Nat.ofDigits} on a cons gives the displayed identity.

  Everything below in this section is derived from this one step, which keeps
  the induction arguments free of any reasoning about digit \emph{lists}.
\end{proof}

\begin{lemma}[Digits of the enumeration]
  \label{lem:digits-enum}
  \lean{IntegerCantorSet.digits_enum}
  \leanok
  \uses{def:enum, lem:digitOf-ne-zero, lem:digitOf-mem, def:standing}
  Assume the standing assumption. Then
  \[
    \dgt_b(\enum(n)) = \bigl(\dgt_\delta(n)\bigr) \text{ with each entry }
      c \text{ replaced by } d_c .
  \]
\end{lemma}
\begin{proof}
  \uses{lem:digits-enum}
  \leanok
  Apply \texttt{Nat.digits\_ofDigits}. Its two hypotheses are: every entry of
  the substituted list is $< b$, which holds since $D \subseteq
  \{0,\ldots,b-1\}$ and Lemma \ref{lem:digitOf-mem}; and the last entry is
  non-zero, which holds by Lemma \ref{lem:digitOf-ne-zero} because the last
  base $\delta$ digit of $n$ is non-zero.
\end{proof}

\begin{lemma}
  \label{lem:enum-mem}
  \lean{IntegerCantorSet.enum_mem_cantorSet}
  \leanok
  \uses{lem:digits-enum, def:cantorSet}
  $\enum(n) \in \Cbd$ for every $n$.
\end{lemma}
\begin{proof}
  \uses{lem:enum-mem}
  \leanok
  Immediate from Lemma \ref{lem:digits-enum} and Lemma \ref{lem:digitOf-mem}.
\end{proof}

\begin{lemma}
  \label{lem:enum-strictMono}
  \lean{IntegerCantorSet.enum_strictMono}
  \leanok
  \uses{lem:enum-recursion, lem:digitOf-strictMono}
  $\enum$ is strictly monotone.
\end{lemma}
\begin{proof}
  \uses{lem:enum-strictMono}
  \leanok
  Strong induction on $m$, using the single-digit recursion
  (Lemma \ref{lem:enum-recursion}) to write
  $\enum(n) = b\,\enum(\lfloor n/\delta \rfloor) + d_{n \bmod \delta}$
  and likewise for $m$. If $\lfloor n/\delta \rfloor <
  \lfloor m/\delta \rfloor$ the induction hypothesis gives
  $\enum(\lfloor n/\delta \rfloor) + 1 \leq
  \enum(\lfloor m/\delta \rfloor)$, and $d_{n \bmod \delta} < b$ absorbs
  the low digit. If the quotients agree then $n \bmod \delta <
  m \bmod \delta$ and Lemma \ref{lem:digitOf-strictMono} applies. This is
  the usual lexicographic comparison, organized so that no statement about
  digit \emph{lists} is needed.
\end{proof}

\begin{theorem}[Proposition 2.5(i)]
  \label{thm:enum-bij}
  \lean{IntegerCantorSet.range_enum}
  \leanok
  \uses{lem:enum-mem, lem:enum-strictMono}
  Assume the standing assumption. Then $\enum$ is an order isomorphism from
  $\N$ onto $\Cbd$; in particular $\enum$ is the increasing enumeration
  $k_0 < k_1 < \cdots$ of $\Cbd$, and $\enum(0) = 0$.
\end{theorem}
\begin{proof}
  \uses{thm:enum-bij}
  \leanok
  Injectivity and order preservation are Lemma \ref{lem:enum-strictMono}.
  Surjectivity onto $\Cbd$: strong induction on $m \in \Cbd$. For $m > 0$ the
  lowest base $b$ digit $m \bmod b$ lies in $D$, so it is $d_c$ for a (unique,
  by Lemma \ref{lem:digitOf-strictMono}) $c < \delta$, and
  $\lfloor m/b \rfloor$ again lies in $\Cbd$; if
  $\enum(k) = \lfloor m/b \rfloor$ then $\enum(\delta k + c) = m$ by
  Lemma \ref{lem:enum-recursion}.  The Lean statement is the range equality
  \texttt{Set.range (enum b D) = cantorSet b D}; combined with
  Lemma \ref{lem:enum-strictMono} this is the order isomorphism, and
  $\enum(0) = 0$ is \texttt{enum\_zero}.
\end{proof}

\begin{lemma}[Proposition 2.5(ii)]
  \label{lem:enum-block}
  \lean{IntegerCantorSet.enum_block}
  \leanok
  \uses{lem:enum-recursion}
  For $i \geq 1$, $m \geq 0$ and $0 \leq j < \delta^i$,
  \[
    \enum(\delta^i m + j) = b^i \enum(m) + \enum(j) .
  \]
\end{lemma}
\begin{proof}
  \uses{lem:enum-block}
  \leanok
  The base $\delta$ digits of $\delta^i m + j$ in positions $0, \ldots, i-1$
  are those of $j$ (padded with zeros), and in positions $\geq i$ are those of
  $m$.  Formally this is an induction on $i$ from the single-digit recursion
  (Lemma \ref{lem:enum-recursion}), writing
  $\delta^{i+1} m + j = \delta(\delta^i m + \lfloor j/\delta\rfloor) +
  j \bmod \delta$; the hypothesis $i \geq 1$ is not needed, since $i = 0$
  forces $j = 0$.
\end{proof}

\begin{lemma}
  \label{lem:sumDigits-enum-block}
  \lean{IntegerCantorSet.sumDigits_enum_block}
  \leanok
  \uses{lem:enum-block, lem:growth, def:sumDigits}
  With $i, m, j$ as above,
  $\sd(\enum(\delta^i m + j)) = \sd(\enum(m)) + \sd(\enum(j))$.
\end{lemma}
\begin{proof}
  \uses{lem:sumDigits-enum-block}
  \leanok
  By Lemma \ref{lem:enum-block} the number in question is
  $b^i \enum(m) + \enum(j)$ with $\enum(j) < b^i$ by Lemma
  \ref{lem:growth}, so the two digit lists concatenate and the sums add.
\end{proof}

\begin{lemma}[Initial segments]
  \label{lem:enum-initial}
  \lean{IntegerCantorSet.enum_image_Iio}
  \leanok
  \uses{thm:enum-bij, lem:growth}
  For every $\ell$, $\enum$ maps $\{0, \ldots, \delta^\ell - 1\}$ bijectively
  onto $\Cbd \cap [0, b^\ell)$. In words, the first $\delta^\ell$ elements of
  the Cantor set are exactly its elements below $b^\ell$.
\end{lemma}
\begin{proof}
  \uses{lem:enum-initial}
  \leanok
  One inclusion is Lemma \ref{lem:enum-mem} together with
  Lemma \ref{lem:growth}.  Conversely, if $k \in \Cbd$ then $k = \enum(n)$
  for some $n$ by Theorem \ref{thm:enum-bij}, and if $n \geq \delta^\ell$
  then $\enum(n) \geq \enum(\delta^\ell) = b^\ell d_1 \geq b^\ell$ by
  monotonicity, so $k < b^\ell$ forces $n < \delta^\ell$.
\end{proof}

\section{Divisibility}

\begin{lemma}
  \label{lem:stepSize-dvd}
  \lean{IntegerCantorSet.stepSize_dvd,
    IntegerCantorSet.stepSize_dvd_sumDigits}
  \leanok
  \uses{def:stepSize, def:cantorSet, def:sumDigits}
  $s(D)$ divides every element of $\Cbd$ and $s(D)$ divides $\sd(k)$ for every
  $k \in \Cbd$.
\end{lemma}
\begin{proof}
  \uses{lem:stepSize-dvd}
  \leanok
  Every digit of $k$ is a member of $D$, hence divisible by $\gcd D$; now sum,
  either weighted by the powers of $b$ (giving $s(D) \mid k$) or unweighted
  (giving $s(D) \mid \sd(k)$).
\end{proof}

\begin{lemma}
  \label{lem:b-sub-one-dvd}
  \lean{IntegerCantorSet.modEq_sumDigits}
  \leanok
  \uses{def:sumDigits}
  $b - 1$ divides $k - \sd(k)$ for every $k$, in the sense that
  $k \equiv \sd(k) \pmod{b-1}$.
\end{lemma}
\begin{proof}
  \uses{lem:b-sub-one-dvd}
  \leanok
  This is \texttt{Nat.modEq\_digits\_sum} applied with $b' = b$ and $b = b-1$,
  using $b \bmod (b-1) = 1$.
\end{proof}

\chapter{Intersective and van der Corput properties}
\label{chap:vdc}

This chapter formalizes Section 3 of the paper, excluding the subsection on
polynomial recurrence, which is Chapter \ref{chap:recurrence}. It is
self-contained: no measure theory and no infinite processes beyond elementary
limits.

\section{Definitions}

\begin{definition}[Avoiding sets]
  \label{def:maxAvoiding}
  \lean{IntegerCantorSet.maxAvoiding}
  \leanok
  For $H \subseteq \N$ and $N \in \N$, let
  \[
    \Lambda(H, N) = \max \{\abs{A} : A \subseteq [N], \ (A - A) \cap
      (H \setminus \{0\}) = \emptyset \} .
  \]
\end{definition}

\begin{definition}[Intersectivity constant]
  \label{def:intersectiveConst}
  \lean{IntegerCantorSet.intersectiveConst}
  \leanok
  \uses{def:maxAvoiding}
  $I(H,N) = \Lambda(H,N)/N$.
\end{definition}

\begin{definition}[Intersective]
  \label{def:isIntersective}
  \lean{IntegerCantorSet.IsIntersective}
  \leanok
  \uses{def:intersectiveConst}
  $H$ is \emph{intersective} if $I(H,N) \to 0$ as $N \to \infty$.
\end{definition}

\begin{lemma}[Subadditivity of $\Lambda$]
  \label{lem:intersective-converges}
  \lean{IntegerCantorSet.maxAvoiding_subadditive}
  \leanok
  \uses{def:maxAvoiding}
  $\Lambda(H, M + N) \leq \Lambda(H, M) + \Lambda(H, N)$ for all $M, N$.
\end{lemma}
\begin{proof}
  \uses{lem:intersective-converges}
  \leanok
  Let $A \subseteq [M+N]$ be avoiding, and split it as $A_1 = A \cap [M]$ and
  $A_2 = A \setminus [M]$. Then $A_1 \subseteq [M]$ is avoiding, and
  $A_2 - M \subseteq [N]$ is avoiding because a translation does not change
  differences; so
  $\abs{A} = \abs{A_1} + \abs{A_2} \leq \Lambda(H,M) + \Lambda(H,N)$.

  (Fekete's lemma, \texttt{Subadditive.tendsto\_lim} in Mathlib, then gives
  that $I(H,N) = \Lambda(H,N)/N$ converges. That consequence is not itself
  formalized: every bound below is proved for $I(H,N)$ directly, so the
  existence of the limit is never used.)
\end{proof}

\begin{definition}[van der Corput constant]
  \label{def:vdcConst}
  \lean{IntegerCantorSet.vdcConst}
  \leanok
  $\tau(H,N)$ is the infimum of $a_0$ over all real sequences
  $(a_n)_{n < N}$ supported on $H \setminus \{0\}$ for which the cosine
  polynomial
  \[
    \Psi(x) = a_0 + \sum_{0 < n < N,\ n \in H} a_n \cos(2\pi n x)
  \]
  satisfies $\Psi \geq 0$ everywhere and $\Psi(0) = 1$.
  Note the infimum is over a non-empty set, since $\Psi \equiv 1$ is
  admissible, so $\tau(H,N) \leq 1$; and $\tau(H,N) \geq 0$ because
  $\int_0^1 \Psi = a_0$.
\end{definition}

\begin{definition}[van der Corput set]
  \label{def:isVdC}
  \lean{IntegerCantorSet.IsVdC}
  \leanok
  \uses{def:vdcConst}
  $H$ is \emph{van der Corput} if $\tau(H,N) \to 0$.
\end{definition}

\begin{lemma}
  \label{lem:vdcConst-antitone}
  \lean{IntegerCantorSet.vdcConst_antitone}
  \leanok
  \uses{def:vdcConst}
  $N \mapsto \tau(H, N)$ is antitone.
\end{lemma}
\begin{proof}
  \uses{lem:vdcConst-antitone}
  \leanok
  Enlarging $N$ enlarges the set of admissible $\Psi$, so the infimum can only
  decrease.
\end{proof}

\section{Two elementary principles}

The next two lemmas are the whole content of the quantitative results below.
Both say that a large ``difference set'' inside $H$ forces $H$ to be
intersective, respectively van der Corput, with an explicit rate.

\begin{lemma}[Translate packing]
  \label{lem:packing}
  \lean{IntegerCantorSet.card_le_of_differences_mem}
  \leanok
  \uses{def:maxAvoiding}
  Let $E \subseteq \Z$ be finite and non-empty, and $K \subseteq \N$, and
  suppose every non-zero difference of two elements of $E$ has absolute value
  in $K$. If $A \subseteq [N]$ satisfies $(A - A) \cap K = \emptyset$, then
  the translates $A + e$, $e \in E$, are pairwise disjoint and contained in an
  interval of length $N + \max E - \min E$, whence
  \[
    \abs{A} \leq \frac{N + \max E - \min E}{\abs{E}} .
  \]
\end{lemma}
\begin{proof}
  \uses{lem:packing}
  \leanok
  If $a + e = a' + e'$ with $e \neq e'$ then $a - a' = e' - e$ is a non-zero
  difference of elements of $E$, so $\abs{a - a'} \in K$ and
  $(A-A) \cap K \neq \emptyset$. Hence $(a,e) \mapsto a + e$ is injective on
  $A \times E$, and it maps into the interval
  $[1 + \min E, N + \max E]$; the formalized statement is the resulting
  $\abs{A}\abs{E} \leq N + \max E - \min E$.
\end{proof}

\begin{lemma}[Fej\'er kernel of a finite set]
  \label{lem:fejer}
  \lean{IntegerCantorSet.vdcConst_le_of_differences_mem}
  \leanok
  \uses{def:vdcConst}
  Let $u_0 < u_1 < \cdots < u_{n-1}$ be naturals all of whose differences
  $u_j - u_i$, $i < j$, lie in $H \setminus \{0\}$ and are $< N$. Then
  $\tau(H, N) \leq 1/n$.
\end{lemma}
\begin{proof}
  \uses{lem:fejer}
  \leanok
  Take $\Psi(x) = \abs{n^{-1}\sum_{j<n} e(u_j x)}^2$. Then $\Psi \geq 0$ and
  $\Psi(0) = 1$. Expanding, $\Psi(x) = n^{-2}\sum_{i,j} e((u_j - u_i)x)$; the
  terms $(i,j)$ and $(j,i)$ combine into a real cosine, so $\Psi$ is a cosine
  polynomial with frequencies in $H \setminus \{0\}$ below $N$, and the $n$
  diagonal terms give $a_0 = 1/n$.

  In Lean the delicate point is that $\Psi$ must be exhibited in the exact
  shape required by Definition \ref{def:vdcConst}, i.e. as a coefficient
  function $a : \N \to \R$. It is convenient to count \emph{ordered} pairs:
  with $d(i,j) = \abs{u_i - u_j}$ put
  $a_m = n^{-2}\abs{\{(i,j) \in [n]^2 : d(i,j) = m\}}$ for every $m$, which
  for $m = 0$ counts the $n$ diagonal pairs and so gives $a_0 = 1/n$
  automatically. Summing over the fibres of $d$ turns $\Psi$ into
  $n^{-2}\sum_{i,j} \cos(\theta_i - \theta_j)$
  with $\theta_i = 2\pi u_i x$, and $\cos(\theta_i - \theta_j) =
  \cos\theta_i\cos\theta_j + \sin\theta_i\sin\theta_j$ gives
  \[
    \Psi(x) = n^{-2}\Bigl[\bigl(\textstyle\sum_i \cos\theta_i\bigr)^2
      + \bigl(\textstyle\sum_i \sin\theta_i\bigr)^2\Bigr] ,
  \]
  which is manifestly $\geq 0$ and equals $1$ at $x = 0$. This stays inside
  $\R$, avoiding \texttt{Complex.normSq} altogether.
\end{proof}

\section{Cantor sets are van der Corput, with a $1/\log N$ rate}

This is Proposition 3.3 of the paper. It is the natural first target of the
whole formalization: elementary, quantitative, and short.

\begin{definition}[The tower]
  \label{def:tower}
  \lean{IntegerCantorSet.tower}
  \leanok
  \uses{def:standing}
  Fix $d \in D \setminus \{0\}$. For $j \geq 0$ set
  $u_j = \sum_{t < j} d\, b^t$, so that $u_0 = 0$.
\end{definition}

\begin{lemma}
  \label{lem:tower-diff}
  \lean{IntegerCantorSet.tower_sub_mem}
  \leanok
  \uses{def:tower, def:cantorSet}
  For $i < j$, $u_j - u_i = \sum_{i \leq t < j} d\, b^t \in \Cbd \setminus
  \{0\}$.
\end{lemma}
\begin{proof}
  \uses{lem:tower-diff}
  \leanok
  The base $b$ digits of $u_j - u_i$ are $d$ in positions $i, \ldots, j-1$ and
  $0$ elsewhere, both of which lie in $D$. (The pair $(i,j)$ is in fact
  recovered from the digit string, but Lemma \ref{lem:fejer} does not need it,
  so this is not formalized.)
\end{proof}

\begin{lemma}
  \label{lem:tower-lt}
  \lean{IntegerCantorSet.tower_lt}
  \leanok
  \uses{def:tower}
  $u_{n-1} = d\,\frac{b^{n-1} - 1}{b-1} \leq b^{n-1} - 1$, so all the
  differences above are $< b^{n-1}$.
\end{lemma}
\begin{proof}
  \uses{lem:tower-lt}
  \leanok
  A geometric sum, using $d \leq b-1$.
\end{proof}

\begin{theorem}[Proposition 3.3]
  \label{thm:logSavings}
  \lean{IntegerCantorSet.vdcConst_cantorSet_le}
  \leanok
  \uses{lem:fejer, lem:tower-diff, lem:tower-lt, lem:vdcConst-antitone}
  Assume the standing assumption. For all integers $N \geq 2$,
  \[
    \tau(\Cbd, N) \leq \frac{\log b}{\log N} .
  \]
  In particular $\Cbd$ is a van der Corput set.
\end{theorem}
\begin{proof}
  \uses{thm:logSavings}
  \leanok
  By Lemmas \ref{lem:tower-diff} and \ref{lem:tower-lt} the hypotheses of
  Lemma \ref{lem:fejer} hold with $N_0 = b^{n-1}$, giving
  $\tau(\Cbd, b^{n-1}) \leq 1/n$. Given $N \geq 2$ take
  $n = 1 + \lfloor \log N / \log b \rfloor$, so $b^{n-1} \leq N$ and
  $n \geq \log N / \log b$; conclude by Lemma \ref{lem:vdcConst-antitone}.
\end{proof}

\begin{theorem}
  \label{thm:logSavings-intersective}
  \lean{IntegerCantorSet.intersectiveConst_cantorSet_le,
    IntegerCantorSet.isIntersective_cantorSet}
  \leanok
  \uses{lem:packing, lem:tower-diff, lem:tower-lt}
  Assume the standing assumption. Then
  $I(\Cbd, N) \leq 2\log b / \log N$ for $N \geq b$, and $\Cbd$ is
  intersective.
\end{theorem}
\begin{proof}
  \uses{thm:logSavings-intersective}
  \leanok
  Take $n = 1 + \lfloor \log N / \log b \rfloor$ and apply
  Lemma \ref{lem:packing} with $E = \{u_0, \ldots, u_{n-1}\}$ and
  $K = \Cbd \setminus \{0\}$, so $\abs{E} = n$ and
  $\max E - \min E < b^{n-1} \leq N$, giving
  $\Lambda(\Cbd, N) \leq 2N/n$ and hence $I(\Cbd,N) \leq 2/n \leq
  2 \log b / \log N$. Intersectivity follows by squeezing, as for
  Theorem \ref{thm:logSavings}.

  We deliberately route around the paper's $I(H,N) \ll \tau(H,N)$, which is a
  theorem of Ruzsa and would be a substantial separate formalization. The
  direct argument gives the same bound up to the constant.
\end{proof}

\begin{lemma}[Necessity of $0 \in D$]
  \label{lem:zero-necessary}
  \lean{IntegerCantorSet.exists_dvd_of_isIntersective}
  \leanok
  \uses{def:isIntersective, def:cantorSet}
  If $H \subseteq \N$ is intersective then for every $q \geq 1$, $H$ contains
  a non-zero multiple of $q$.
\end{lemma}
\begin{proof}
  \uses{lem:zero-necessary}
  \leanok
  If $H$ contains no non-zero multiple of $q$, take $A = q\Z \cap [N]$; then
  $A - A \subseteq q\Z$ misses $H$. Formalized with $N = qM$, where
  $A = \{q, 2q, \ldots, Mq\}$ has exactly $M$ elements, so
  $I(H, qM) \geq 1/q$ for every $M \geq 1$ and $I(H,N) \not\to 0$.

  (The consequence drawn in the paper --- if $0 \notin D$ then $\Cbd$ is not
  intersective, because every element of $\Cbd$ then has non-zero last digit
  and so is not divisible by $b$ --- is not formalized; the standing
  assumption includes $0 \in D$ throughout.)
\end{proof}

\section{Signed digits and power savings}

This section formalizes \S3.1. The digit set is now allowed to contain
negative integers, so $\Cbd \subseteq \Z$ and the definitions of
Section \ref{chap:vdc} are applied to $H = \abs{\Cbd} \setminus \{0\}$.

\begin{definition}[Signed Cantor set]
  \label{def:signedCantorSet}
  \lean{IntegerCantorSet.signedCantorSet}
  Let $R_b \subseteq (-b, b)$ be a complete residue system mod $b$ with
  $0 \in R_b$, and $D \subseteq R_b$ with $0 \in D$ and
  $2 \leq \abs{D} < b$. Define $\Cbd \subseteq \Z$ by the same digit formula.
  Uniqueness of the digit representation still holds, by induction on the
  least significant digit.
\end{definition}

\begin{theorem}[Proposition 3.7]
  \label{thm:powerSavings}
  \lean{IntegerCantorSet.vdcConst_signed_le}
  \uses{def:signedCantorSet, def:vdcConst}
  Let $S$ be a finite set of integers with $S - S \subseteq D$ and
  $\abs{S} = \sigma > 1$. Then for $N \geq b$,
  \[
    \tau(\Cbd, N) \leq \sigma N^{-\log \sigma / \log b} .
  \]
\end{theorem}
\begin{proof}
  \uses{thm:powerSavings}
  Take $P_J(x) = \prod_{j \leq J} \abs{\sigma^{-1}\sum_{s \in S}
  e(s b^j x)}^2$. Expanding, the frequencies are $\sum_j (s_j - t_j) b^j$,
  whose digits $s_j - t_j$ lie in $S - S \subseteq D \subseteq R_b
  \subseteq (-b, b)$, so the frequencies lie in $\Cbd$ and are
  $< b^{J+1}$ in absolute value. By uniqueness of digits only the
  $\sigma^{J+1}$ diagonal terms contribute to the constant term, giving
  $a_0 = \sigma^{-(J+1)}$. Interpolate using
  Lemma \ref{lem:vdcConst-antitone}.
\end{proof}

\chapter{The key exponential sum}
\label{chap:expsum}

This chapter formalizes Lemma 4.7 of the paper, the single result on which
Chapters \ref{chap:pudu}, \ref{chap:residue} and \ref{chap:ergodic} all rest.
It is unblocked but long.  It is now formalized, in
\texttt{IntegerCantorSets/ExpSum.lean}.

Two pieces of notation are given Lean definitions of their own here:
$e(x) = \exp(2\pi i x)$ is \texttt{IntegerCantorSet.expc}, and
$\norm{x}$, the distance from $x$ to the nearest integer, is
\texttt{IntegerCantorSet.distInt}, defined as $\abs{x - \mathrm{round}(x)}$.
The quantitative bridge between them, $\abs{1 - e(x)} = 2\sin(\pi\norm{x})
\geq 4\norm{x}$, is \texttt{four\_mul\_distInt\_le}.

\begin{definition}
  \label{def:Lhat}
  \lean{IntegerCantorSet.Lhat}
  \leanok
  $\Lhat(\theta) = \delta^{-1} \sum_{d \in D} e(d\theta)$.
\end{definition}

\begin{lemma}
  \label{lem:Lhat-abs-one}
  \lean{IntegerCantorSet.abs_Lhat_eq_one_iff}
  \leanok
  \uses{def:Lhat, def:stepSize}
  Assume $0 \in D$. Then $\abs{\Lhat(\theta)} = 1$ if and only if
  $s(D)\,\theta \in \Z$. In particular $\abs{\Lhat(\theta)} < 1$ except at
  the $s(D)$ points of $[0,1)$ with $s(D)\theta \in \Z$.
\end{lemma}
\begin{proof}
  \uses{lem:Lhat-abs-one}
  \leanok
  Equality in the triangle inequality for a mean of unit vectors forces all
  the $e(d\theta)$ to be equal; since $0 \in D$ that common value is $1$, i.e.
  $d\theta \in \Z$ for all $d \in D$, which is equivalent to
  $\gcd(D)\,\theta \in \Z$ by B\'ezout.

  Formalized through the quantitative defect estimate
  \texttt{norm\_Lhat\_le}: for every $d \in D$,
  \[
    \abs{\Lhat(\theta)} \leq 1 - \frac{\abs{1 - e(d\theta)}^2}{4\delta},
  \]
  which comes from isolating the two terms $d = 0$ and $d$ in the sum and
  using $\abs{1 + z}^2 + \abs{1 - z}^2 = 4$ for $\abs{z} = 1$.  Equality
  $\abs{\Lhat(\theta)} = 1$ therefore forces $e(d\theta) = 1$ for every
  $d \in D$.  The B\'ezout step is \texttt{exists\_gcd\_eq\_sum}, proved by
  induction on the finite set: $\gcd(D) = \sum_{d \in D} f_d\, d$ for some
  integers $f_d$.  The same estimate is reused in
  Theorem \ref{thm:irrational}, where the convergence
  $\abs{\Lhat(\theta_c)} \to 1$ is turned into $\norm{d\theta_c} \to 0$
  without any compactness argument.
\end{proof}

\begin{lemma}[Block product]
  \label{lem:blockProduct}
  \lean{IntegerCantorSet.blockProduct}
  \leanok
  \uses{lem:enum-recursion, def:Lhat}
  With $\Phi(n) = e(\enum(n)\alpha + \sd(\enum(n))\beta)$,
  \[
    P_\ell := \delta^{-\ell} \sum_{0 \leq j < \delta^\ell} \Phi(j)
      = \prod_{c=0}^{\ell-1} \Lhat(b^c \alpha + \beta) .
  \]
\end{lemma}
\begin{proof}
  \uses{lem:blockProduct}
  \leanok
  The digits of $\enum(j)$ range independently and uniformly over $D$ as $j$
  ranges over $[0, \delta^\ell)$, and
  $\enum(j)\alpha + \sd(\enum(j))\beta = \sum_c d_{j_c}(b^c\alpha + \beta)$.
  This is a product of independent averages.

  Formalized by induction on $\ell$ with $\alpha$ generalized, peeling off the
  \emph{least} significant base $\delta$ digit rather than a whole block: with
  $n = \delta m + j$ and $j < \delta$, Lemma \ref{lem:enum-recursion} and
  \texttt{sumDigits\_mul\_add} give
  $\Phi_{\alpha,\beta}(n) = \Phi_{b\alpha,\beta}(m)\, e(d_j(\alpha + \beta))$,
  so $P_\ell(\alpha) = \Lhat(\alpha + \beta)\, P_{\ell-1}(b\alpha)$, which is
  the recursion satisfied by the product (\texttt{Finset.prod\_range\_succ'}).
  That the inner sum over $j < \delta$ is $\delta\,\Lhat(\alpha + \beta)$ is
  \texttt{sum\_range\_card\_digitOf}, the statement that $\texttt{digitOf}$ is a
  bijection $[0,\delta) \to D$.
\end{proof}

\begin{lemma}[Window splitting]
  \label{lem:windowSplit}
  \lean{IntegerCantorSet.window_split}
  \leanok
  \uses{lem:enum-block, lem:blockProduct}
  For $0 \leq M \leq N$ and $i \geq 0$, with
  $M_i = \lceil M/\delta^i \rceil$ and $N_i = \lfloor N/\delta^i \rfloor$, put
  $\Psi_i(m) = e(b^i \enum(m)\alpha + \sd(\enum(m))\beta)$.  Then
  \[
    \abs{\sum_{M \leq n < N} \Phi(n)
      - \delta^i P_i \sum_{M_i \leq m < N_i} \Psi_i(m)} \leq 2\delta^i .
  \]
  (The Lean statement is this un-normalised inequality rather than the
  normalised identity-with-$O$ of the paper; dividing by $N - M$ and using
  $\delta^i(N_i - M_i) \leq N - M$ recovers the latter.  The hypothesis
  $i \geq 1$ is not needed: $i = 0$ is trivial.)
\end{lemma}
\begin{proof}
  \uses{lem:windowSplit}
  \leanok
  Lemma \ref{lem:enum-block} factors $\Phi$ over each complete block
  $[\delta^i m, \delta^i(m+1))$, giving $\Psi_i(m) \cdot \delta^i P_i$ there
  by Lemma \ref{lem:blockProduct}; at most $2\delta^i$ indices lie outside
  the complete blocks, and each $\abs{\Phi(n)} = 1$.  In Lean the two
  elementary facts $M \leq \delta^i M_i < M + \delta^i$ and
  $\delta^i N_i \leq N < \delta^i N_i + \delta^i$ are isolated as
  \texttt{ceil\_div\_bounds} and \texttt{floor\_div\_bounds}; the degenerate
  case $N_i < M_i$ (no complete block) is handled separately, where they give
  $N - M < \delta^i$ directly.  The consequence
  $\delta^i(N_i - M_i) \leq N - M$ is \texttt{block\_count\_le}, and the
  normalised window bound actually used in Theorem \ref{thm:irrational} is
  \texttt{norm\_window\_le}.
\end{proof}

\begin{definition}[Exceptional pairs]
  \label{def:exceptional}
  \lean{IntegerCantorSet.Exceptional}
  \leanok
  $(\alpha, \beta)$ is \emph{exceptional} if there are integers $a$, $r$ and
  $t \geq 0$ with
  \[
    s\alpha \equiv \frac{a}{b-1} + \frac{r}{b^t}, \qquad
    s\beta \equiv -\frac{a}{b-1} \pmod 1 .
  \]
  (That the set of exceptional pairs is countable is clear from the
  definition and is not needed below, so it is not formalized.)
\end{definition}

\begin{lemma}
  \label{lem:escape}
  \lean{IntegerCantorSet.exceptional_of_tendsto}
  \leanok
  \uses{def:exceptional}
  If $\norm{s(b^c\alpha + \beta)} \to 0$ as $c \to \infty$ then
  $(\alpha, \beta)$ is exceptional, and moreover there is a $t$ with
  $s(b^c\alpha + \beta) \in \Z$ for all $c \geq t$.
\end{lemma}
\begin{proof}
  \uses{lem:escape}
  \leanok
  Write $\alpha' = s\alpha$, $\beta' = s\beta$ and put $x = (b-1)\alpha'$, so
  $\norm{b^c x} \leq \norm{b^{c+1}\alpha' + \beta'} +
  \norm{b^c \alpha' + \beta'} \to 0$. Choose $c_0$ with
  $\norm{b^c x} < 1/(2b)$ for $c \geq c_0$. If $\norm{b^{c_0}x} \neq 0$ then
  $\norm{b^{c_0+i}x} = b^i\norm{b^{c_0}x}$ eventually exceeds $1/(2b)$, a
  contradiction. Hence $b^{c_0}x \in \Z$, and partial fractions using
  $\gcd(b-1, b^t) = 1$ give the stated form with $t = c_0$.  Finally, for
  $c \geq t$ the quantity $b^c\alpha' + \beta'$ differs from
  $a/(b-1) + \beta'$ by an integer, so $\norm{a/(b-1) + \beta'}$ is both
  constant in $c$ and a limit of a null sequence, hence zero.

  The $D$-free core is \texttt{escape\_core}; the amplification step
  $\norm{by} = b\norm{y}$ when $\norm{y} < 1/(2b)$ is
  \texttt{distInt\_nat\_mul}.
\end{proof}

\begin{theorem}[Lemma 4.7]
  \label{thm:irrational}
  \lean{IntegerCantorSet.tendsto_expSum}
  \leanok
  \uses{lem:Lhat-abs-one, lem:blockProduct, lem:windowSplit, lem:escape}
  The limit
  $\lim_{N} N^{-1}\!\!\sum_{0 \leq n < N} \Phi(n)$
  exists for all $(\alpha, \beta)$, and is zero unless $(\alpha,\beta)$ is
  exceptional. If $(\alpha,\beta)$ is not exceptional, the convergence to
  zero is uniform over windows: the average over $M \leq n < N$ tends to zero
  as $N - M \to \infty$.
\end{theorem}
\begin{proof}
  \uses{thm:irrational}
  \leanok
  Write $\theta_c = b^c\alpha + \beta$ and $P_\ell = \prod_{c<\ell}
  \Lhat(\theta_c)$.  Each $\abs{\Lhat(\theta_c)} \leq 1$, so
  $\abs{P_\ell}$ is non-increasing.  Split on whether it tends to $0$.

  If $\abs{P_\ell} \to 0$, feeding $\abs{P_i} < \varepsilon/2$ into Lemma
  \ref{lem:windowSplit} gives $\abs{\sum_{M \leq n < N}\Phi(n)} \leq
  (\varepsilon/2)(N-M) + 2\delta^i$, hence the uniform window statement, and
  with $M = 0$ the limit $0$.

  Otherwise $\abs{P_\ell} \geq \varepsilon_0 > 0$ for all $\ell$.  Being
  non-increasing and bounded below, $\abs{P_\ell}$ converges to a positive
  limit, so $\abs{\Lhat(\theta_c)} = \abs{P_{c+1}}/\abs{P_c} \to 1$.  The
  defect estimate of Lemma \ref{lem:Lhat-abs-one} then gives
  $\abs{1 - e(d\theta_c)} \to 0$, hence $\norm{d\theta_c} \to 0$, for every
  $d \in D$; B\'ezout turns this into $\norm{s\theta_c} \to 0$.  Lemma
  \ref{lem:escape} applies: $(\alpha,\beta)$ is exceptional and
  $s\theta_c \in \Z$ for $c \geq t$.  Since $s \mid d$ for every $d \in D$,
  every factor with $c \geq t$ equals $1$, so $P_\ell = P_t$ for
  $\ell \geq t$, and the inner sum in Lemma \ref{lem:windowSplit} at $i = t$
  is identically $1$: for $n \in \Cbd$ with digits $d_j$,
  $b^t n \alpha + \sd(n)\beta = \sum_j d_j \theta_{t+j} \in \Z$
  (\texttt{exists\_int\_of\_mem\_cantorSet}, by induction on $n$ through
  $n = b\lfloor n/b\rfloor + (n \bmod b)$).  Comparing $\delta^t\lfloor
  N/\delta^t\rfloor$ with $N$ costs a further $\delta^t$, so the limit is
  $P_t$.

  Two deviations from the paper's route.  The reduction to $s(D) = 1$ is not
  carried out: $s$ is kept explicit throughout, which is why Lemma
  \ref{lem:escape} and Definition \ref{def:exceptional} carry the factor $s$.
  And the first case is settled by monotonicity of $\abs{P_\ell}$ together
  with the quantitative defect estimate, rather than by compactness; this
  also makes Lemma \ref{lem:b-sub-one-dvd} unnecessary here, its role being
  taken by the digit decomposition just described.
\end{proof}

\chapter{Polynomial uniform distribution}
\label{chap:pudu}

This chapter formalizes Section 4; the file is
\texttt{IntegerCantorSets/PUDU.lean}. Everything below is formalized, with no
\texttt{sorry}: Theorem \ref{thm:polyUniform} is proved unconditionally.

Polynomials in two variables over $\R$ are represented by
\texttt{MvPolynomial (Fin 2) $\R$}, evaluated at integer points
(\texttt{polyEval}). The blueprint's ring $(\R/\Q)[X,Y]$ is not formed: a
quotient ring is unnecessary, because ``$\overline{Q}$ vanishes'' is simply
``every coefficient is rational'', and Mathlib has \texttt{Irrational} on
individual reals. This keeps every statement of the chapter first order in the
coefficients.

\begin{definition}
  \label{def:irrationalPoly}
  \lean{IntegerCantorSet.IsIrrationalPoly, IntegerCantorSet.IrrDegLE}
  \leanok
  A polynomial $p \in \R[X,Y]$ is \emph{irrational} if some non-constant
  monomial has an irrational coefficient (\texttt{IsIrrationalPoly}). Its
  \emph{irrational degree is at most} $\ell$ if every monomial of total degree
  greater than $\ell$ has a rational coefficient (\texttt{IrrDegLE}).

  The irrational degree is not introduced as a number. Every statement below is
  phrased with the ``at most'' predicate, which is what the induction needs,
  and \texttt{irrDegLE\_totalDegree} supplies the bound $\ell = \deg p$, so the
  predicate is always available for some $\ell$.
\end{definition}

\begin{definition}[PUDU]
  \label{def:PUDU}
  \lean{IntegerCantorSet.IsPUDU, IntegerCantorSet.UniformAvg}
  \leanok
  \uses{def:irrationalPoly}
  A sequence $(x_n)$ in $\Z^2$ is \emph{PUDU} if for every irrational
  polynomial $p$ and every $\varepsilon > 0$ there is a $K$ such that
  \[
    \norm{ \frac{1}{N-M}\sum_{M \leq n < N} \e{p(x_n)} } < \varepsilon
    \qquad \text{whenever } M + K \leq N .
  \]
  This is exactly the uniform-window shape of the conclusion of Theorem
  \ref{thm:irrational}, reused rather than restated.
  \texttt{UniformAvg} names the displayed condition for a single $p$, and
  \texttt{cantorSeq} is the sequence $n \mapsto (\enum(n), \sd(\enum(n)))$.
\end{definition}

\begin{lemma}[Bad directions]
  \label{lem:badDirections}
  \lean{IntegerCantorSet.badDirections}
  \leanok
  \uses{def:irrationalPoly}
  Let $p \in \R[X,Y]$ have an irrational coefficient on some monomial of total
  degree $\ell \geq 1$. Then there are integers $(c_1,c_2) \neq (0,0)$ such
  that every $w \in \Z^2$ for which $\partial_w p$ has only rational
  coefficients in degree $\ell - 1$ satisfies $c_1w_1 + c_2w_2 = 0$.

  Here $\partial_w p = w_1 \partial_X p + w_2 \partial_Y p$
  (\texttt{dirDeriv}, built from Mathlib's \texttt{MvPolynomial.pderiv}); the
  blueprint's condition $\partial_w \overline{Q} = 0$ is the rationality of
  those coefficients. Directions are taken in $\Z^2$, not $\Q^2$, because that
  is all Theorem \ref{thm:UDstep} supplies.
\end{lemma}
\begin{proof}
  \uses{lem:badDirections}
  \leanok
  Let $m_0$ be the monomial of degree $\ell$ with irrational coefficient. Since
  $\ell \geq 1$, $m_0$ involves $X$ or $Y$; say $m_0 = m + e_X$ with
  $\deg m = \ell - 1$ (the other case is symmetric). Put
  \[
    u = (m_X + 1)\,p_{m + e_X}, \qquad v = (m_Y + 1)\,p_{m + e_Y},
  \]
  so that $u$ is irrational. By \texttt{MvPolynomial.coeff\_pderiv} the
  coefficient of $m$ in $\partial_w p$ is $w_1 u + w_2 v$, so it suffices to
  show (\texttt{exists\_bad\_direction}) that the $w \in \Z^2$ with
  $w_1u + w_2v \in \Q$ lie on a rational line, given that $u$ or $v$ is
  irrational.

  If the only such $w$ is $0$, take $(c_1,c_2) = (1,0)$. Otherwise fix
  $w^0 \neq 0$ with $w^0_1u + w^0_2v \in \Q$ and take
  $(c_1,c_2) = (-w^0_2, w^0_1)$, which is non-zero. For any other such $w$ set
  $d = w^0_1w_2 - w^0_2w_1$. Then
  \[
    w^0_1(w_1u + w_2v) - w_1(w^0_1u + w^0_2v) = d\,v, \qquad
    w^0_2(w_1u + w_2v) - w_2(w^0_1u + w^0_2v) = -d\,u ,
  \]
  and both left-hand sides are rational. If $d \neq 0$ this makes $u$ and $v$
  both rational, a contradiction; so $d = 0$, which is the assertion
  $c_1w_1 + c_2w_2 = 0$. No quotient ring and no dimension theory over $\Q$ is
  used.
\end{proof}

\begin{theorem}[Base case]
  \label{thm:UD1}
  \lean{IntegerCantorSet.UD_one}
  \leanok
  \uses{thm:irrational, def:PUDU}
  Assume the standing assumption. The conclusion of Definition \ref{def:PUDU}
  holds for the sequence $n \mapsto (\enum(n), \sd(\enum(n)))$ and every
  irrational $p$ of irrational degree at most $1$.
\end{theorem}
\begin{proof}
  \uses{thm:UD1}
  \leanok
  Let $\alpha$ and $\beta$ be the coefficients of $X$ and $Y$ in $p$, and let
  $g = p - p_0 - \alpha X - \beta Y$ (\texttt{ratPart}). Since $p$ has
  irrational degree at most $1$ and is irrational, $\alpha$ or $\beta$ is
  irrational, and every coefficient of $g$ is rational.

  A product of the denominators over the support of $g$ gives a $q \geq 1$
  with $q\,g_m \in \Z$ for every $m$ (\texttt{exists\_common\_denom}); hence
  $g(x,y) - g(x',y') \in \Z$ whenever $x \equiv x'$ and $y \equiv y' \pmod q$
  (\texttt{polyEval\_sub\_int\_of\_dvd}), so $\e{g(x,y)}$ depends only on
  $(x,y)$ modulo $q$. Expanding the indicator of each residue class by
  orthogonality of the additive characters of $\Z/q$,
  \[
    \frac{1}{q}\sum_{t < q} \e{\frac{tz}{q}}
      = \begin{cases} 1 & q \mid z \\ 0 & \text{otherwise,}\end{cases}
  \]
  (\texttt{sum\_expc\_orthogonality}, a geometric sum) and reordering the
  resulting four-fold sum gives the finite Fourier expansion
  \[
    \e{g(x,y)} = \sum_{s,t < q} c_{s,t}\, \e{\frac{sx + ty}{q}}
  \]
  (\texttt{fourier\_expansion}). Therefore
  \[
    \e{p(\enum(n), \sd(\enum(n)))}
      = \sum_{s,t<q} \e{p_0} c_{s,t}\,
        \Phi_{\alpha + s/q,\ \beta + t/q}(n) ,
  \]
  a fixed finite linear combination of the sums of Theorem
  \ref{thm:irrational}. Each shifted pair $(\alpha + s/q, \beta + t/q)$ has an
  irrational coordinate, hence is not exceptional: an exceptional pair has
  $s(D)\alpha$ and $s(D)\beta$ both rational by Definition
  \ref{def:exceptional} (\texttt{not\_exceptional\_of\_irrational}, which uses
  $s(D) \geq 1$, \texttt{stepSize\_pos}). So Theorem \ref{thm:irrational}
  applies to each of the $q^2$ terms; taking $K$ to be the largest of the $q^2$
  thresholds and $\varepsilon/(1 + \sum_{s,t}\abs{c_{s,t}})$ as the tolerance
  for each gives the claim.
\end{proof}

\begin{theorem}[Induction step]
  \label{thm:UDstep}
  \lean{IntegerCantorSet.UD_succ}
  \leanok
  \uses{lem:badDirections, lem:enum-block, lem:Lhat-abs-one, def:PUDU}
  Assume the standing assumption and let $\ell \geq 1$. If the conclusion of
  Definition \ref{def:PUDU} holds for every irrational polynomial of
  irrational degree at most $\ell$, then it holds for every irrational
  polynomial of irrational degree at most $\ell + 1$.
\end{theorem}
\begin{proof}
  \uses{thm:UDstep}
  \leanok
  We may assume $p$ has a monomial of total degree exactly $\ell+1$ with
  irrational coefficient, since otherwise the inductive hypothesis applies
  directly. Let $(c_1,c_2)$ be the direction supplied by Lemma
  \ref{lem:badDirections}.

  Fix $\rho$ and write $\Lambda = \delta^\rho$. Splitting $[M,N)$ into the
  complete blocks of length $\Lambda$ it contains costs at most $2\Lambda$
  unimodular terms (\texttt{window\_block\_split}, the same ceiling/floor
  bookkeeping as Lemma \ref{lem:windowSplit}). On a block, Lemma
  \ref{lem:enum-block} gives
  $\enum(\Lambda m + j) = b^\rho \enum(m) + \enum(j)$ and
  $\sd(\enum(\Lambda m + j)) = \sd(\enum(m)) + \sd(\enum(j))$. Cauchy--Schwarz
  in the block index (\texttt{norm\_sum\_sq\_le\_card\_mul}) and the identity
  \[
    \abs{\sum_{j<\Lambda} \e{t_j}}^2
      = \sum_{j,j'<\Lambda} \e{t_j - t_{j'}}
  \]
  (\texttt{normSq\_sum\_expc}) replace $p$ by the \emph{difference
  polynomials}
  \[
    R_{j,j'}(X,Y) = p\bigl(b^\rho X + \enum(j),\ Y + \sd(\enum(j))\bigr)
      - p\bigl(b^\rho X + \enum(j'),\ Y + \sd(\enum(j'))\bigr)
  \]
  (\texttt{blockDiff}), which satisfy
  $R_{j,j'}(\enum(m), \sd(\enum(m))) = p(x_{\Lambda m + j}) - p(x_{\Lambda m+j'})$.

  Writing the coefficients of an affine substitution as a single sum over the
  support (\texttt{coeff\_substAffine}) shows two things at once. First, the
  degree $\ell+1$ coefficients of $R_{j,j'}$ cancel and the higher ones are
  rational, so $R_{j,j'}$ has irrational degree at most $\ell$
  (\texttt{irrDegLE\_diffPoly}). Second, in degree $\ell$ only the three
  monomials $m$, $m + e_X$, $m + e_Y$ contribute, and the result is
  $b^{\rho m_X} \partial_w p$ up to a rational, where
  $w = (\enum(j) - \enum(j'), \sd(\enum(j)) - \sd(\enum(j')))$
  (\texttt{coeff\_diffPoly\_sub\_isRat}). Hence, by Lemma
  \ref{lem:badDirections}, $R_{j,j'}$ is irrational --- so the inductive
  hypothesis applies to it --- unless $(j,j')$ is \emph{bad}, meaning
  $c_1\enum(j) + c_2\sd(\enum(j))$ agrees for $j$ and $j'$.

  The proportion of bad pairs tends to $0$. If $c_1 \neq 0$ then a bad partner
  of $j$ is determined by its digit sum, because $\enum$ is injective; since
  $\enum(j') < b^\rho$ forces $\sd(\enum(j')) \leq (b-1)\rho$, there are at
  most $(b-1)\rho + 1$ of them, and $((b-1)\rho+1)/\delta^\rho \to 0$. If
  $c_1 = 0$ then bad means $\sd(\enum(j)) = \sd(\enum(j'))$, and
  orthogonality of the characters $\theta \mapsto \e{k\theta}$ on $[0,1]$
  turns the count into
  \[
    \#\{(j,j') : \sd(\enum(j)) = \sd(\enum(j'))\}
      = \int_0^1 \abs{\sum_{j<\Lambda} \e{\sd(\enum(j))\theta}}^2 d\theta
      = \delta^{2\rho}\int_0^1 \abs{\Lhat(\theta)}^{2\rho}\,d\theta ,
  \]
  the middle step by Lemma \ref{lem:blockProduct} with $\alpha = 0$. By Lemma
  \ref{lem:Lhat-abs-one}, $\abs{\Lhat(\theta)} = 1$ only when
  $s(D)\theta \in \Z$, a countable set; so $\abs{\Lhat(\theta)}^{2\rho} \to 0$
  almost everywhere and the integral tends to $0$ by dominated convergence.

  Choosing $\rho$ with bad proportion below $\varepsilon^2/8$, and the
  inductive tolerance also $\varepsilon^2/8$, makes the square of the
  normalised block average at most $\varepsilon^2/4$; the incomplete blocks
  contribute $2\Lambda/(N-M)$, which is below $\varepsilon/2$ once $N-M$ is
  large.
\end{proof}

\begin{theorem}[Theorem 4.3]
  \label{thm:polyUniform}
  \lean{IntegerCantorSet.isPUDU, IntegerCantorSet.UD_le}
  \leanok
  \uses{thm:UD1, thm:UDstep}
  Assume the standing assumption. Then
  $\bigl((\enum(n), \sd(\enum(n)))\bigr)_{n \geq 0}$ is PUDU.
\end{theorem}
\begin{proof}
  \uses{thm:polyUniform}
  \leanok
  Induction on $\ell$ in ``irrational degree at most $\ell$''
  (\texttt{UD\_le}): the base case is Theorem \ref{thm:UD1} and the step is
  Theorem \ref{thm:UDstep}. Any irrational $p$ has irrational degree at most
  $\deg p$ (\texttt{irrDegLE\_totalDegree}), which supplies the $\ell$.
\end{proof}

\chapter{Distribution in residue classes}
\label{chap:residue}

This chapter formalizes Section 5; the file is
\texttt{IntegerCantorSets/Residue.lean}. Everything below is formalized, with
no \texttt{sorry}. It is finite Fourier analysis on $\Z_q$ layered over
Theorem \ref{thm:irrational}, which supplies every Fourier coefficient that
occurs; the polynomial uniform distribution of Chapter \ref{chap:pudu} is not
needed here.

Counting is done with
\[
  \texttt{countMod}\,f\,q\,x\,N = \#\{n < N : f(n) \equiv x \bmod q\},
  \qquad
  \texttt{densityMod}\,f\,q\,x\,N = \tfrac1N \texttt{countMod}\,f\,q\,x\,N ,
\]
and \texttt{countMod}$_2$, \texttt{densityMod}$_2$ are the two-dimensional
versions, for a pair of sequences and a pair of moduli. ``The sequence has a
limiting distribution $\pi_q$'' is formalized as: for each $x$ separately, the
real sequence $N \mapsto \texttt{densityMod}\,f\,q\,x\,N$ converges. Since
$\Z_q$ is finite this is the same as weak convergence of measures, and no
measure theory is needed. Residues are naturals throughout, so $x$ ranges over
$\N$ and the statements are invariant under $x \mapsto x + q$ by construction.

\begin{lemma}[Fourier inversion on $\Z_q$]
  \label{lem:fourierMod}
  \lean{IntegerCantorSet.countMod_fourier, IntegerCantorSet.tendsto_densityMod,
    IntegerCantorSet.countMod2_fourier, IntegerCantorSet.tendsto_densityMod2}
  \leanok
  For $q \geq 1$ and any $f : \N \to \N$,
  \[
    \texttt{countMod}\,f\,q\,x\,N
      = \frac1q \sum_{u < q} \e{-\tfrac{ux}{q}} \sum_{n<N} \e{\tfrac{u f(n)}{q}} .
  \]
  Consequently, if for every $u < q$ the average
  $N^{-1}\sum_{n<N} \e{u f(n)/q}$ converges to $c(u)$, then
  $\texttt{densityMod}\,f\,q\,x\,N$ converges to
  $\frac1q\sum_{u<q} \e{-ux/q}\,c(u)$, for every $x$. The same two statements
  hold on $\Z_a \times \Z_{a'}$ for a pair $(f,g)$.
\end{lemma}
\begin{proof}
  \uses{lem:fourierMod}
  \leanok
  Orthogonality of the additive characters of $\Z_q$
  (\texttt{sum\_expc\_orthogonality}, already available from Chapter
  \ref{chap:pudu}) gives
  $\frac1q\sum_{u<q}\e{u(m-x)/q} = \mathbf 1_{m \equiv x \bmod q}$; summing
  this over $n < N$ at $m = f(n)$ and exchanging the two sums is the displayed
  identity. For the limit, divide by $N$ and apply the identity with the inner
  average in place of the inner sum; the outer sum is finite, so the limit
  passes through it. The two-dimensional statement multiplies the two
  indicators and exchanges three sums.

  The limit is produced in $\C$ and converted to $\R$ through continuity of
  $\mathrm{Re}$: \texttt{tendsto\_densityMod} takes the real value as a
  hypothesis (``this Fourier inverse equals $v$''), while
  \texttt{tendsto\_densityMod\_re} states the limit as the real part of the
  inverse, which is what the existence statements below use.
\end{proof}

\begin{proposition}[Proposition 5.1]
  \label{prop:sumDigits}
  \lean{IntegerCantorSet.sumDigits_equidistributed,
    IntegerCantorSet.sumDigits_uniform_iff}
  \leanok
  \uses{thm:irrational, def:stepSize, lem:fourierMod}
  With $g = \gcd(q, s)$, the sequence $\sd(\enum(n)) \bmod q$ is
  asymptotically uniformly distributed on the subgroup $g\Z_q$: for every $x$,
  \[
    \lim_{N \to \infty} \frac{\#\{n<N : \sd(\enum(n)) \equiv x \bmod q\}}{N}
      = \begin{cases} g/q & g \mid x \\ 0 & \text{otherwise.}\end{cases}
  \]
  In particular it is asymptotically uniformly distributed on all of $\Z_q$ if
  and only if $\gcd(q,s) = 1$.
\end{proposition}
\begin{proof}
  \uses{prop:sumDigits}
  \leanok
  Every digit is a multiple of $s$, so $g \mid s \mid \sd(\enum(n))$ for all
  $n$ (\texttt{stepSize\_dvd\_sumDigits}). The Fourier coefficients of the
  uniform measure on $g\Z_q$ are $\mathbf 1_{(q/g) \mid u}$, so by Lemma
  \ref{lem:fourierMod} it suffices to show
  \[
    \lim_{N\to\infty} \frac1N \sum_{n<N} \e{\sd(\enum(n))\tfrac uq}
      = \mathbf 1_{(q/g)\mid u}, \qquad u < q .
  \]
  If $(q/g) \mid u$ every summand is $1$, because $g \mid \sd(\enum(n))$.
  Otherwise apply Theorem \ref{thm:irrational} at $(\alpha,\beta) = (0,u/q)$:
  the limit is $0$ unless $(0,u/q)$ is exceptional. If it were, then
  $0 = \frac{a}{b-1} + \frac{r}{b^t} + k$ for integers $a,r,k$; clearing
  denominators and using $\gcd(b-1,b^t) = 1$ forces $(b-1) \mid a$
  (\texttt{sub\_one\_dvd\_of\_eq\_zero}), so $\frac{a}{b-1} \in \Z$ and then
  $s\frac uq \in \Z$, i.e.\ $q \mid su$. Since
  $\gcd(q/g, s/g) = 1$ this gives $(q/g) \mid u$
  (\texttt{div\_gcd\_dvd\_of\_dvd\_mul}), a contradiction.

  The Fourier inverse is evaluated once and for all in
  \texttt{fourier\_subgroup}: reindexing $u = (q/g)v$ turns
  $\frac1q\sum_{u<q}\e{-ux/q}\mathbf 1_{(q/g)\mid u}$ into
  $\frac1q\sum_{v<g}\e{-vx/g}$, which orthogonality evaluates to $g/q$ or $0$.
  The ``in particular'' is the case $x = 1$ together with uniqueness of
  limits.
\end{proof}

\begin{lemma}[Lemma 5.3]
  \label{lem:Allq}
  \lean{IntegerCantorSet.limitDistribution,
    IntegerCantorSet.limitDistribution_coprime,
    IntegerCantorSet.limitDistribution_uniform_iff,
    IntegerCantorSet.limitDistribution_dvd_pow,
    IntegerCantorSet.limitDistribution_mul}
  \leanok
  \uses{thm:irrational, lem:enum-block, lem:fourierMod}
  For every $q \geq 1$, $\enum(n) \bmod q$ has a limiting distribution $\pi_q$
  on $\Z_q$. Moreover:
  \begin{enumerate}
    \item If $\gcd(q,b) = 1$ then $\pi_q$ is uniform on $g\Z_q$,
      $g = \gcd(q,s)$; it is uniform on $\Z_q$ if and only if $\gcd(q,s)=1$.
    \item If $q \mid b^j$ then $\pi_q(x)$ is
      $\delta^{-j}\#\{i < \delta^j : \enum(i) \equiv x \bmod q\}$, the
      distribution of $\Cbd \cap [0,b^j) \bmod q$.
    \item If $q = ru$ with $\gcd(r,b) = 1$ and $u \mid b^j$, then
      $\pi_q(x) = \pi_r(x)\,\pi_u(x)$ for every $x$, with $\pi_r$ as in (1) and
      $\pi_u$ as in (2).
  \end{enumerate}
  Part (3) is the product measure $\pi_r \times \pi_u$ of the paper: $\gcd(r,u)
  = 1$ is automatic from $\gcd(r,b)=1$ and $u \mid b^j$, and under the Chinese
  remainder isomorphism $\Z_q \cong \Z_r \times \Z_u$ the pair attached to $x$
  is $(x \bmod r, x \bmod u)$, so the displayed identity for all $x$ is exactly
  the product statement. The isomorphism itself is never formed.
\end{lemma}
\begin{proof}
  \uses{lem:Allq}
  \leanok
  Existence: by Theorem \ref{thm:irrational} the limit of
  $\frac1N\sum_{n<N}\e{\enum(n)\frac uq}$ exists for every $u$, and Lemma
  \ref{lem:fourierMod} turns this into convergence of each $\pi_q(x)$.

  (1) As in Proposition \ref{prop:sumDigits}, with $s \mid \enum(n)$ in place
  of $s \mid \sd(\enum(n))$ and $(\alpha,\beta) = (u/q,0)$. If $(u/q,0)$ were
  exceptional then $0 \equiv -\frac{c}{b-1} \bmod 1$ puts $\frac{c}{b-1}$ in
  $\Z$ directly, whence $b^t\frac{su}{q} \in \Z$; as $\gcd(q,b)=1$ this gives
  $q \mid su$ and again $(q/g) \mid u$.

  (2) By Lemma \ref{lem:enum-block}, for $i < \delta^j$ one has
  $\enum(\delta^jm+i) = b^j\enum(m) + \enum(i) \equiv \enum(i) \bmod q$, so a
  complete block of length $\delta^j$ contributes exactly
  $A = \#\{i<\delta^j : \enum(i) \equiv x\}$ (\texttt{countMod\_block}), and
  $\texttt{countMod}(\delta^jM) = MA$. For general $N$, write
  $N = \delta^jK + F$ with $F < \delta^j$; monotonicity of the count and the
  bound $C(N) \leq C(M) + (N-M)$ for $M \leq N$, where $C$ is the count, give
  \[
    \abs{ \tfrac{1}{N}C(N) - A\delta^{-j} } \leq \delta^j/N ,
  \]
  which tends to $0$. This replaces the paper's $O(\delta^j/N)$ by an explicit
  inequality.

  (3) Let $\Lambda = \delta^j$ and $w = b^j$, so $\gcd(w,r) = 1$. For
  $i < \Lambda$ let $y_i < r$ be the unique solution of
  $w\,y + \enum(i) \equiv x \bmod r$ (\texttt{linSol}, defined through
  $(w)^{-1}$ in \texttt{ZMod r}, and characterized by \texttt{lin\_cond\_iff}).
  Since $\enum(\Lambda m + i) = w\,\enum(m) + \enum(i)$ and $u \mid w$, the
  Chinese remainder theorem gives
  \[
    \enum(\Lambda m+i) \equiv x \bmod ru
    \iff \enum(i) \equiv x \bmod u \ \text{ and }\ \enum(m) \equiv y_i \bmod r,
  \]
  so, summing over a block,
  $\texttt{countMod}\,q\,x\,(\Lambda M)
    = \sum_{i \in S} \texttt{countMod}\,r\,y_i\,M$ with
  $S = \{i<\Lambda : \enum(i) \equiv x \bmod u\}$. Dividing by $\Lambda M$ and
  letting $M \to \infty$, part (1) evaluates each inner limit as $\pi_r(y_i)$.
  Finally $g' = \gcd(r,s)$ divides $\enum(i)$, so
  $w\,y_i \equiv x \bmod g'$, and $\gcd(g',b) = 1$ gives
  $g' \mid y_i \iff g' \mid x$; hence $\pi_r(y_i) = \pi_r(x)$ for every $i$ and
  the limit is $\abs{S}\Lambda^{-1}\pi_r(x) = \pi_u(x)\pi_r(x)$. The limit over
  all $N$ exists by the first paragraph, so it agrees with this limit over the
  subsequence $N = \Lambda M$; no second error analysis is needed.
\end{proof}

Note in particular that if $q \mid b^j$ for some $j \geq 1$, the Cantor set
could be uniformly distributed on $\Z_q$, depending on the choice of $D$ and
the smallest $j$ with $q \mid b^j$.

\begin{lemma}[Lemma 5.5]
  \label{lem:positiveDensity}
  \lean{IntegerCantorSet.positive_density}
  \leanok
  \uses{lem:Allq}
  For every $q \geq 1$, the set $\{n : q \mid \enum(n)\}$ has a density, and
  that density is strictly positive.
\end{lemma}
\begin{proof}
  \uses{lem:positiveDensity}
  \leanok
  The density exists by Lemma \ref{lem:Allq} at $x = 0$; call it $c$. For
  positivity the formalization does not go through
  $\pi_q(0) = \pi_r(0)\pi_u(0)$, so part (3) of Lemma \ref{lem:Allq} is not
  used: a one-sided bound along a subsequence is enough and is shorter.

  Take $j = q$, $g = \gcd(q,b^q)$ and $r = q/g$. Then $\gcd(r,b)=1$
  (\texttt{coprime\_div\_gcd\_pow}: if a prime $p$ divided both, then with
  $e = v_p(q) < q \le j$ one has $p^e \mid g$ and $p \mid r$, so
  $p^{e+1} \mid q$, contradicting $e = v_p(q)$). Since $g \mid b^q$ and
  $q = gr$,
  \[
    r \mid \enum(m) \implies q \mid b^q\enum(m) = \enum(\delta^q m),
  \]
  so $m \mapsto \delta^q m$ injects $\{m < K : r \mid \enum(m)\}$ into
  $\{n < \delta^qK : q \mid \enum(n)\}$. Dividing by $\delta^qK$ and letting
  $K \to \infty$ along the subsequence $N = \delta^qK$,
  \[
    c \;\geq\; \delta^{-q}\,\frac{\gcd(r,s)}{r} \;>\; 0 ,
  \]
  the inner factor by part (1) of Lemma \ref{lem:Allq} applied to $r$ at
  $x = 0$.
\end{proof}

\begin{theorem}[Theorem 5.7]
  \label{thm:modularJoint}
  \lean{IntegerCantorSet.joint_limitDistribution, IntegerCantorSet.JointUniform,
    IntegerCantorSet.joint_uniform,
    IntegerCantorSet.not_jointUniform_of_gcd_enum,
    IntegerCantorSet.not_jointUniform_of_gcd_sumDigits,
    IntegerCantorSet.not_jointUniform_of_gcd_sub_one}
  \leanok
  \uses{lem:Allq, prop:sumDigits, lem:fourierMod}
  Let $a, a' \geq 1$ and consider the pair
  $(\enum(n) \bmod a,\ \sd(\enum(n)) \bmod a') \in \Z_a \times \Z_{a'}$.
  \begin{enumerate}
    \item The limiting distribution exists.
    \item It is uniform (the predicate \texttt{JointUniform}: every class has
      density $(aa')^{-1}$) if $\gcd(a,sb) = 1$ and $\gcd(a',s(b-1)) = 1$.
    \item It is not uniform if $\gcd(a,s) > 1$, or $\gcd(a',s) > 1$, or
      $\gcd(a,a',b-1) > 1$.
  \end{enumerate}
  The sufficient condition in (2) is not necessary; with $b=4$, $D=\{0,1\}$,
  $\enum(n) \bmod 2$ is uniform although $\gcd(2,b)>1$, and
  $\sd(\enum(n)) \bmod 3$ is uniform by Proposition \ref{prop:sumDigits}
  although $\gcd(3,b-1)>1$. (That remark is not formalized.)
\end{theorem}
\begin{proof}
  \uses{thm:modularJoint}
  \leanok
  (1) By Theorem \ref{thm:irrational} the limit of
  $\frac1N\sum_{n<N}\e{\enum(n)\frac ua + \sd(\enum(n))\frac{u'}{a'}}$ exists
  for all $u,u'$; the two-dimensional half of Lemma \ref{lem:fourierMod} turns
  this into convergence of each class.

  (2) By Lemma \ref{lem:fourierMod} it suffices that the limit above vanish for
  $u < a$, $u' < a'$ not both $0$, and the $(0,0)$ term is $1$; the Fourier
  inverse then collapses to $(aa')^{-1}$. By Theorem \ref{thm:irrational} it is
  enough that $(u/a, u'/a')$ not be exceptional. If it were, then for integers
  $c, r, t$,
  \[
    s\frac ua = \frac{c}{b-1} + \frac{r}{b^t}, \qquad
    s\frac{u'}{a'} + \frac{c}{b-1} \in \Z .
  \]
  Put $Y = c/(b-1)$. Both $a'Y$ and $(b-1)Y$ are integers, and
  $\gcd(a',b-1)=1$, so B\'ezout makes $Y$ an integer
  (\texttt{isInt\_of\_coprime}). Hence $a' \mid su'$, and $\gcd(a',s)=1$ gives
  $a' \mid u'$, so $u' = 0$. Then $b^tsu/a \in \Z$, and $\gcd(a,sb)=1$ gives
  $a \mid u$, so $u = 0$: a contradiction.

  (3) In each case a single residue class is empty for every $N$, which
  contradicts uniformity because $(aa')^{-1} > 0$
  (\texttt{not\_jointUniform\_of\_countMod2\_eq\_zero}). If
  $d = \gcd(a,s) > 1$, the class $(x,y) = (1,0)$ is empty: $d \mid \enum(n)$
  for all $n$, so $\enum(n) \equiv 1 \bmod a$ would force $d \mid 1$. If
  $d = \gcd(a',s) > 1$, the class $(0,1)$ is empty by the same argument applied
  to $\sd(\enum(n))$. If $d = \gcd(a,\gcd(a',b-1)) > 1$, the class $(0,1)$ is
  empty because $\enum(n) \equiv \sd(\enum(n)) \bmod (b-1)$
  (Lemma \ref{lem:b-sub-one-dvd}), so modulo $d$ one would get $0 \equiv 1$.
\end{proof}

\chapter{Polynomial recurrence}
\label{chap:recurrence}

\begin{definition}[Upper Banach density]
  \label{def:banach}
  \lean{IntegerCantorSet.upperBanachDensity}
  \notready
  $d^*(A) = \lim_{L \to \infty} \sup_{t} \abs{A \cap [t, t+L)}/L$.
  Not currently in Mathlib.
\end{definition}

\begin{theorem}[Furstenberg correspondence]
  \label{thm:furstenberg}
  \notready
  \uses{def:banach}
  For $A \subseteq \Z$ with $d^*(A) > 0$ there is an invertible measure
  preserving system $(X, \mathcal{B}, \mu, T)$ and $\tilde A \subseteq X$ with
  $\mu(\tilde A) = d^*(A)$ and
  $d^*(A \cap (A-m)) \geq \mu(\tilde A \cap T^{-m}\tilde A)$ for all $m$.

  \textbf{Not in Mathlib.} This is a prerequisite project in its own right.
  Until it lands, Theorem \ref{thm:polysIntersective} should be stated in its
  dynamical form, which is what the paper's proof actually establishes.
\end{theorem}

\begin{theorem}[Theorem 3.14, dynamical form]
  \label{thm:polysIntersective}
  \lean{IntegerCantorSet.liminf_recurrence}
  \notready
  \uses{thm:polyUniform, lem:positiveDensity, thm:furstenberg}
  Let $p$ be integer valued with no constant term, and let
  $(X,\mathcal{B},\mu,T)$ be an invertible measure preserving system with
  $\mu(\tilde A) > 0$. Then
  \[
    \liminf_N \frac1N \sum_{0 \leq n < N}
      \mu\bigl(\tilde A \cap T^{-p(\enum(n))}\tilde A\bigr) > 0 .
  \]
  Deducing that $H_p$ is intersective additionally needs the weak$^*$
  compactness argument of the paper, which is formalizable once
  Definition \ref{def:banach} exists.
\end{theorem}

\chapter{Ergodic theorems}
\label{chap:ergodic}

\begin{theorem}[Theorem 6.1]
  \label{thm:ergodic}
  \notready
  \uses{thm:irrational}
  For commuting unitaries $(U,V)$ on a Hilbert space and all $x$,
  $\frac1N\sum_{0 \leq n < N} U^{\enum(n)}V^{\sd(\enum(n))}x$ converges in
  norm, to $\sum_{(\alpha,\beta) \in \mathcal{E}} F(\alpha,\beta)
  P_{\alpha,\beta}x$.

  \textbf{Blocked.} The proof is a two-line consequence of Theorem
  \ref{thm:irrational} \emph{given} the joint spectral measure $E$ of
  $(U,V)$ on $\T^2$. Mathlib does not have it. A worthwhile intermediate
  target is the scalar statement: the averages
  $\frac1N\sum e(\enum(n)\alpha + \sd(\enum(n))\beta)$ converge, which is
  exactly Theorem \ref{thm:irrational} and is already in scope.
\end{theorem}

\begin{theorem}[Theorem 6.4]
  \label{thm:polyMeanErgodic}
  \notready
  \uses{thm:polyUniform, thm:modularJoint, thm:ergodic}
  The polynomial version, for an integer valued $p : \Z^2 \to \Z^2$.
\end{theorem}

\chapter{Metric pair correlation}
\label{chap:paircorr}

\begin{definition}
  \label{def:pairCorrelation}
  \lean{IntegerCantorSet.HasPairCorrelation}
  A sequence $(\theta_n)$ has \emph{Poissonian pair correlation} if for every
  $s > 0$,
  $N^{-1}\abs{\{0 \leq j \neq k < N : \norm{\theta_j - \theta_k} < s/N\}}
  \to 2s$.
\end{definition}

\begin{definition}[Additive energy]
  \label{def:additiveEnergy}
  \lean{IntegerCantorSet.additiveEnergy}
  $E(A) = \abs{\{(x,y,z,w) \in A^4 : x + y = z + w\}}$.
  Mathlib has \texttt{Finset.additiveEnergy}; check whether the conventions
  agree before reproving anything.
\end{definition}

\begin{lemma}[No carries]
  \label{lem:noCarry}
  \lean{IntegerCantorSet.additiveEnergy_pow}
  \uses{def:additiveEnergy, lem:enum-initial}
  If $D \subseteq \{0, \ldots, \lfloor (b-1)/2 \rfloor\}$ then addition of two
  elements of $\Cbd \cap [0,b^\ell)$ produces no carries, and consequently
  $E(\Cbd \cap [0,b^\ell)) = E(D)^\ell$.
\end{lemma}
\begin{proof}
  \uses{lem:noCarry}
  Each digit of the sum is $x_j + y_j \leq b-1$, so the base $b$ digits of
  $x+y$ are the digitwise sums; hence $x + y = z + w$ iff
  $x_j + y_j = z_j + w_j$ for all $j$, and the count factors.

  This lemma is entirely finitary and is a good second target, independent of
  everything in Chapters \ref{chap:expsum}--\ref{chap:residue}.
\end{proof}

\begin{lemma}
  \label{lem:growth}
  \lean{IntegerCantorSet.enum_lt_pow}
  \leanok
  \uses{lem:enum-recursion}
  If $n < \delta^\ell$ then $\enum(n) < b^\ell$.
\end{lemma}
\begin{proof}
  \uses{lem:growth}
  \leanok
  Induction on $\ell$ from the single-digit recursion
  (Lemma \ref{lem:enum-recursion}): $\enum(n) = b\,
  \enum(\lfloor n/\delta \rfloor) + d_{n \bmod \delta}$ with
  $\lfloor n/\delta \rfloor < \delta^{\ell-1}$ and
  $d_{n \bmod \delta} < b$.

  (The polynomial form used in Chapter \ref{chap:paircorr},
  $\enum(n) \leq b\,n^{d}$ with $d = \log b / \log \delta$, follows by
  taking $\ell$ minimal with $n < \delta^{\ell}$, so that
  $\delta^{\ell-1} \leq n$ and $b^{\ell} = b \cdot (\delta^{\ell-1})^{d}
  \leq b\,n^{d}$.  That real-exponent form is not formalized; Chapter
  \ref{chap:paircorr} is blocked on Theorem \ref{thm:ALL} in any case.)
\end{proof}

\begin{theorem}[Aistleitner--Larcher--Lewko]
  \label{thm:ALL}
  \notready
  If $(a(x))$ are distinct non-negative integers with
  $E(\{a(0),\ldots,a(N-1)\}) \ll N^{3-\varepsilon}$, then $(\alpha a(x))$ has
  Poissonian pair correlation for a.e. $\alpha$; if moreover
  $a(x) \ll x^d$, the exceptional set has Hausdorff dimension at most
  $1 - \varepsilon/(d+3)$.

  \textbf{Not a formalization target.} This is recorded as a hypothesis so
  that Theorem \ref{thm:pairCorrelation} can be stated and its reduction
  formalized.
\end{theorem}

\begin{theorem}[Proposition 7.3]
  \label{thm:pairCorrelation}
  \notready
  \uses{lem:noCarry, lem:growth, thm:ALL}
  Assuming Theorem \ref{thm:ALL}, if $D \subseteq
  \{0,\ldots,\lfloor (b-1)/2\rfloor\}$ then for a.e. $\alpha$ the sequence
  $(\alpha\,\enum(n))$ has Poissonian pair correlation, with the stated bound
  on the dimension of the exceptional set.
\end{theorem}
\begin{proof}
  \uses{thm:pairCorrelation}
  Lemmas \ref{lem:noCarry} and \ref{lem:growth} verify the two hypotheses of
  Theorem \ref{thm:ALL}, with $\varepsilon = \varepsilon(D)$ and
  $d = \log b/\log\delta$. The only genuinely new content is the monotonicity
  step $E(A_N) \leq E(D)^\ell \leq \delta^3 N^{3-\varepsilon}$ for general
  $N$, which is Lemma \ref{lem:noCarry} plus monotonicity of additive energy
  under inclusion.
\end{proof}
